\documentclass[12pt,a4paper]{article}

\usepackage[T2A]{fontenc}
\usepackage[utf8]{inputenc}
\usepackage[russian]{babel}
\usepackage{amsmath,amssymb}
\usepackage[margin=2.5cm]{geometry}
\usepackage{enumitem}

\begin{document}

\begin{center}
    {\Large Теоретическое домашнее задание №2}\\[8pt]
\end{center}

\vspace{2em}

\noindent
\textbf{Задача 1.} 

На занятиях мы не успели поговорить о том, как измерить качество модели в задаче классификации. Для этого есть множество интерпретируемых метрик, которые выражаются с помощью \textit{confusion matrix}, в частности: \textit{Accuracy, Precision, Recall, F1-score}; отдельно стоит \textit{ROC-AUC}, который считается по семейству confusion matrix, построенных при всех возможных порогах.


Вам необходимо разобраться в понятиях, выделенных курсивом, выяснить, как они считаются, и дать содержательную интерпретацию. Подумайте, в каких задачах стоит выбирать ту или иную метрику.

\vspace{2em}
\noindent
\textbf{Задача 2.}

Пусть $b_1,\dots,b_M$~--- одинаково распределенные базовые алгоритмы, $a_M(x)=\frac{1}{M}\sum_m b_m(x)$, и в таком случае $\mathbb{E}b_i(x)=\mu$, $Var(b_i(x))=\sigma^2, \,\,\, i=1,\dots,M$
\begin{enumerate}[label=(\alph*)]
	\item Покажите, что $\mathbb{E}\,a_M(x)=\mathbb{E}\,b_1(x)$. Какой вывод отсюда следует о выборе базовой модели для бэггинга?
	\item На занятиях мы пользовались тем фактом, что для $a_M(x)$ дисперсия равна
	
	\[
	\mathrm{Var}(a_M(x))=\rho\sigma^2+\frac{1-\rho}{M}\sigma^2,
	\]
	где $\sigma^2$~--- дисперсия отдельного базового алгоритма, $\rho$~--- корреляция между различными парами базовых алгоритмов $b$. Как получается такая формула? 
	\item Объясните с помощью пунктов (a)--(b), почему бэггинг применяют к глубоким деревьям, а не к линейной регрессии или к пням глубины~1. Что произойдет в формуле из (b), если базовые алгоритмы почти не различаются ($\rho\to1$)?

\end{enumerate}
\vspace{2em}
\noindent
\textbf{Задача 3.} 

Убедитесь, что вы знаете ответы на следующие вопросы:
\begin{itemize}
	\item Как устроен алгоритм построения решающего дерева?
	\item Что выдается в листе дерева в задаче классификации и в задаче регрессии?
	\item В чем идея бэггинга? Какие преимущества его применения?
	\item Как устроен алгоритм обучения случайного леса?
	\item Как kNN применяется в задаче регрессии? Как получить из kNN оценки вероятностей классов?
	\item Как гиперпараметр $k$ влияет на сложность модели? Что происходит при $k=1$ и при $k=n$ (где $n$~--- размер обучающей выборки)?
	\item Нужно ли масштабировать признаки для kNN? А для решающего дерева?
	\item В чем проявляется «проклятие размерности» применительно к kNN?
\end{itemize}


\end{document}
